\subsubsection{Data Persistence}
\begin{enumerate}
    \item \textbf{Array Serialization and Deserialization Module}
          
          The array serialization and deserialization module provides comprehensive routines
          for efficiently storing and retrieving multi-dimensional arrays of various
          data types.
          
          \begin{enumerate}
            \item \textbf{General Architecture}
                  
                  The module supports five data types (integer, real, character, logical
                  and complex) across dimensions 1 through 5. All routines follow a
                  consistent naming convention and parameter structure for ease of use.
                  
            \item \textbf{Supported Functions}
                  
                  \begin{itemize}
                    \item \texttt{serialize\_int\_Nd(arr, filename, ierr)} - Serialize N-dimensional
                          integer arrays.
                          
                    \item \texttt{deserialize\_int\_Nd(arr, filename, ierr)} -
                          Deserialize N-dimensional integer arrays.
                          
                    \item \texttt{serialize\_real\_Nd(arr, filename, ierr)} - Serialize
                          N-dimensional real arrays.
                          
                    \item \texttt{deserialize\_real\_Nd(arr, filename, ierr)} -
                          Deserialize N-dimensional real arrays.
                          
                    \item \texttt{serialize\_char\_Nd(arr, filename, ierr)} - Serialize
                          N-dimensional character arrays.
                          
                    \item \texttt{deserialize\_char\_Nd(arr, filename, ierr)} -
                          Deserialize N-dimensional character arrays.
                          
                    \item \texttt{serialize\_logical\_Nd(arr, filename, ierr)} -
                          Serialize N-dimensional logical arrays.
                          
                    \item \texttt{deserialize\_logical\_Nd(arr, filename, ierr)} -
                          Deserialize N-dimensional logical arrays.
                          
                    \item \texttt{serialize\_complex\_Nd(arr, filename, ierr)} -
                          Serialize N-dimensional complex arrays.
                          
                    \item \texttt{deserialize\_complex\_Nd(arr, filename, ierr)} -
                          Deserialize N-dimensional complex arrays.
                          
                    \item \texttt{get\_array\_metadata(filename, dims, max\_dims, ndims,
                          ierr, [clen])} - Retrieve array metadata from files
                  \end{itemize}
                  
                  Where N represents the array dimension (1-5).
                  
            \item \textbf{Detailed Parameter Specifications}
                  
                  \textbf{Serialization Functions Parameters:}
                  \begin{itemize}
                    \item \texttt{arr} (input): The array to be serialized. \\ \textbf{Type:}
                          Allocatable array of type \texttt{integer(int32)}, \texttt{real(real64)},
                          \texttt{character(len=clen)}, \texttt{(}logical) or \texttt{complex}
                          \\ \textbf{Dimensions:} 1D to 5D \\ \textbf{Intent:} \texttt{in}
                          \\ \textbf{Note:} \textit{clen} has to be the size of the arrays
                          elements
                          
                    \item \texttt{filename} (input): The name of the binary file to create.
                          \\ \textbf{Type:} \texttt{character(len=*)} \\ \textbf{Intent:} \texttt{in}
                          
                    \item \texttt{ierr} (output): Error code indicating operation status.
                          \\ \textbf{Type:} \texttt{integer(int32)} \\ \textbf{Intent:} \texttt{out}
                          \\ \textbf{Values:} 0 indicates success, non-zero values indicate
                          specific error conditions
                  \end{itemize}
                  \begin{lstlisting}[language=Fortran, caption={Example usage of a serialize function}]
call serialize_int_2d(array, "test_file.bin", ierr)
                  \end{lstlisting}
                  
            \item \textbf{Metadata Retrieval Function Parameters:}
                  \begin{itemize}
                    \item \texttt{filename} (input): The name of the binary file to read
                          metadata from. \\ \textbf{Type:} \texttt{character(len=*)} \\
                          \textbf{Intent:} \texttt{in}
                          
                    \item \texttt{dims} (output): Array containing the dimensions of the
                          stored array. \\ \textbf{Type:} \texttt{integer(int32), dimension(max\_dims)}
                          \\ \textbf{Intent:} \texttt{out} \textbf{Note:} Must be preallocated,
                          recommended with size 5
                          
                    \item \texttt{max\_dims} (input): Size of the dims array (typically
                          5 for maximum supported dimensions). \\ \textbf{Type:} \texttt{integer(int32)}
                          \\ \textbf{Intent:} \texttt{in}
                          
                    \item \texttt{ndims} (output): Number of dimensions actually used by
                          the array. \\ \textbf{Type:} \texttt{integer(int32)} \\ \textbf{Intent:}
                          \texttt{out}
                          
                    \item \texttt{ierr} (output): Error code indicating operation status.
                          \\ \textbf{Type:} \texttt{integer(int32)} \\ \textbf{Intent:} \texttt{out}
                          \\ \textbf{Values:} 0 indicates success, non-zero values indicate
                          specific error conditions
                          
                    \item \texttt{clen} (output, optional): For character arrays, the length
                          of character elements. \\ \textbf{Type:} \texttt{integer(int32)}
                          \\ \textbf{Intent:} \texttt{out} \\ \textbf{Note:} This parameter is
                          only used and required for character arrays
                  \end{itemize}
                  \begin{lstlisting}[language=fortran, caption={Usage of metadata retrieval}]
call get_array_metadata(filename, dims, max_dims, ndims, ierr, clen)
                  \end{lstlisting}
                  
                  \textbf{Deserialization Functions Parameters:}
                  \begin{itemize}
                    \item \texttt{arr} (output): The array to be populated with deserialized
                          data. \\ \textbf{Type:} Allocatable array of type \texttt{integer(int32)},
                          \texttt{real(real64)}, \texttt{character(len=clen)}, \texttt{logical},
                          or \texttt{complex} \\ \textbf{Dimensions:} 1D to 5D \\ \textbf{Intent:}
                          \texttt{out} \\ \textbf{Note:} The array must be allocated with
                          correct dimensions and correct element length before calling
                          deserialization routines
                          
                    \item \texttt{filename} (input): The name of the binary file to read.
                          \\ \textbf{Type:} \texttt{character(len=*)} \\ \textbf{Intent:} \texttt{in}
                          
                    \item \texttt{ierr} (output): Error code indicating operation status.
                          \\ \textbf{Type:} \texttt{integer(int32)} \\ \textbf{Intent:} \texttt{out}
                          \\ \textbf{Values:} 0 indicates success, non-zero values indicate
                          specific error conditions
                  \end{itemize}
                  \begin{lstlisting}[language=fortran, caption={Example usage of a deserialize function}]
call deserialize_int_2d(array, "test_file.bin", ierr)
                  \end{lstlisting}
                  \begin{lstlisting}[language=fortran, caption={Full example}]
integer(int32), allocatable :: kallisto_expression(:,:) 
integer(int32), allocatable :: kallisto_expression_2(:,:)
integer(int32) :: dims(5)
integer(int32) :: ndims, ierr, max_dims
character(len=64) :: fname

!Assume array is filled with data

call serialize_real_2d(kallisto_expression, &
                        "kallisto_data_v1.bin", ierr)
if(.not. is_ok(ierr)) error stop ierr

!Reading the data again
call get_array_metadata("kallisto_data_v1.bin", dims, max_dims, &
                        ndims, ierr)
if(.not. is_ok(ierr)) error stop ierr
allocate(kallisto_expression_2(dims(1), dims(2)))

call deserialize_real_2d(kallisto_expression_2, &
                        "kallisto_data_v1.bin", ierr)
if(.not. is_ok(ierr)) error stop ierr
                  \end{lstlisting}
                  
            \item \textbf{Error Handling}
                  
                  All functions return an error code through the \texttt{ierr} parameter.
                  The error handling utility \hyperref[sec:error-handling]{tox errors} provides
                  functions to check and interpret these codes.
                  
            \item \textbf{Implementation Notes}
                  
                  \begin{itemize}
                    \item All routines are implemented in standard Fortran for interoperability
                          
                    \item The module uses allocatable arrays for flexible memory management
                          
                    \item File operations use Fortran's native unformatted I/O for efficiency
                          
                    \item Error handling is consistent across all functions for predictable
                          behavior
                  \end{itemize}
          \end{enumerate}
          
    \item \textbf{Standard Archive Operations}
          \begin{itemize}
            \item \textbf{Save Tox Data Archive}
                  
                  Saves multiple data arrays to a ZIP archive with automatic manifest creation
                  and temporary file management. Handles any combination of available
                  standard data arrays. This function can only be used with standard gene
                  expression data. If other data needs to be saved, see section "Non-standard
                  archive operations". Call the \texttt{save\_tox\_data} subroutine with
                  the following arguments:
                  
                  \textbf{Input arguments:}
                  \begin{itemize}
                    \item \texttt{zip\_filename character(len=*)}: Output ZIP file name
                          
                    \item \texttt{ierr integer(int32)}: Error code
                  \end{itemize}
                  \textbf{Optional input arguments (provide both array and filename for each
                  data type):}
                  \begin{itemize}
                    \item \texttt{gene\_ids character(len=*)}: Gene IDs array
                          
                    \item \texttt{gene\_ids\_file character(len=*)}: Internal filename for
                          gene IDs
                          
                    \item \texttt{expression real(real64)}: Expression data matrix
                          
                    \item \texttt{expression\_file character(len=*)}: Internal filename for
                          expression data
                          
                    \item \texttt{gene\_to\_family integer(int32)}: Gene to family mapping
                          
                    \item \texttt{gene\_to\_family\_file character(len=*)}: Internal filename
                          for mapping
                          
                    \item \texttt{family\_ids character(len=*)}: Family IDs array
                          
                    \item \texttt{family\_ids\_file character(len=*)}: Internal filename
                          for family IDs
                          
                    \item \texttt{family\_centroids real(real64)}: Family centroids matrix
                          
                    \item \texttt{family\_centroids\_file character(len=*)}: Internal filename
                          for centroids
                          
                    \item \texttt{shift\_vectors real(real64)}: Shift vectors matrix
                          
                    \item \texttt{shift\_vectors\_file character(len=*)}: Internal filename
                          for shift vectors
                  \end{itemize}
                  \begin{lstlisting}[language=Fortran, caption={Save Complete Data Archive}]
integer(int32) :: ierr

! Save all available data arrays
call save_tox_data("complete_data.zip", ierr, &
                  gene_ids=gene_ids, &
                  gene_ids_file="genes.bin", &
                  expression=expr_data, &
                  expression_file="expression.bin", &
                  gene_to_family=gene_to_fam, &
                  gene_to_family_file="mapping.bin", &
                  family_ids=family_ids, &
                  family_ids_file="families.bin", &
                  family_centroids=centroids, &
                  family_centroids_file="centroids.bin", &
                  shift_vectors=shift_vecs, &
                  shift_vectors_file="shifts.bin")
                  \end{lstlisting}
                  
                  \begin{lstlisting}[language=Fortran, caption={Save Partial Data Archive}]
! Save only gene IDs and expression data
call save_tox_data("expression_only.zip", ierr, &
                  gene_ids=gene_ids, &
                  gene_ids_file="genes.bin", &
                  expression=expr_data, &
                  expression_file="expression.bin")

! Save only family-related data
call save_tox_data("families_only.zip", ierr, &
                  family_ids=family_ids, &
                  family_ids_file="families.bin", &
                  family_centroids=centroids, &
                  family_centroids_file="centroids.bin")
                  \end{lstlisting}
                  
            \item \textbf{Read Tox Data Archive}
                  
                  Reads data arrays from a ZIP archive created with \texttt{save\_tox\_data}.
                  Automatically handles array allocation based on archive contents. The
                  function will try to read all requested arrays (optional arguments that
                  are provided) from the given archive. If an array can not be found in the
                  archive, the output arrays wont be allocated or returned. Call the \texttt{read\_tox\_data}
                  subroutine with the following arguments:
                  
                  \textbf{Input arguments:}
                  \begin{itemize}
                    \item \texttt{zip\_filename character(len=*)}: Input ZIP file name
                          
                    \item \texttt{ierr integer(int32)}: Error code
                  \end{itemize}
                  \textbf{Optional output arguments:}
                  \begin{itemize}
                    \item \texttt{gene\_ids character(len=:), allocatable}: Gene IDs array
                          
                    \item \texttt{expression real(real64), allocatable}: Expression data
                          matrix
                          
                    \item \texttt{gene\_to\_family integer(int32), allocatable}: Gene to
                          family mapping
                          
                    \item \texttt{family\_ids character(len=:), allocatable}: Family IDs
                          array
                          
                    \item \texttt{family\_centroids real(real64), allocatable}: Family centroids
                          matrix
                          
                    \item \texttt{shift\_vectors real(real64), allocatable}: Shift vectors
                          matrix
                  \end{itemize}
                  \textbf{Optional filename output arguments:}
                  \begin{itemize}
                    \item \texttt{gene\_ids\_file character(len=:), allocatable}: Internal
                          filename used
                          
                    \item \texttt{expression\_file character(len=:), allocatable}: Internal
                          filename used
                          
                    \item \texttt{gene\_to\_family\_file character(len=:), allocatable}:
                          Internal filename used
                          
                    \item \texttt{family\_ids\_file character(len=:), allocatable}: Internal
                          filename used
                          
                    \item \texttt{family\_centroids\_file character(len=:), allocatable}:
                          Internal filename used
                          
                    \item \texttt{shift\_vectors\_file character(len=:), allocatable}: Internal
                          filename used
                  \end{itemize}
                  \begin{lstlisting}[language=Fortran, caption={Read Complete Data Archive}]
character(len=:), allocatable :: loaded_gene_ids(:), loaded_family_ids(:)
real(real64), allocatable :: loaded_expr(:,:), loaded_centroids(:,:)
real(real64), allocatable :: loaded_shifts(:,:)
integer(int32), allocatable :: loaded_gene_to_fam(:)
integer(int32) :: ierr

! Read all available data from archive
call read_tox_data("complete_data.zip", ierr, &
                  gene_ids=loaded_gene_ids, &
                  expression=loaded_expr, &
                  gene_to_family=loaded_gene_to_fam, &
                  family_ids=loaded_family_ids, &
                  family_centroids=loaded_centroids, &
                  shift_vectors=loaded_shifts)
                  \end{lstlisting}
                  
                  \begin{lstlisting}[language=Fortran, caption={Read Partial Data Archive}]
! Read only specific arrays (others remain unallocated)
call read_tox_data("expression_only.zip", ierr, &
                  gene_ids=loaded_gene_ids, &
                  expression=loaded_expr)
                  \end{lstlisting}
                  
                  \begin{lstlisting}[language=Fortran, caption={Read with Filename Retrieval}]
character(len=:), allocatable :: actual_gene_file 
character(len=:), allocatable :: actual_expr_file
integer(int32) :: ierr

! Get both data and internal filenames
call read_tox_data("complete_data.zip", ierr, &
                  gene_ids=loaded_gene_ids, &
                  gene_ids_file=actual_gene_file, &
                  expression=loaded_expr, &
                  expression_file=actual_expr_file)
                  \end{lstlisting}
          \end{itemize}
          
    \item \textbf{Non-Standard Archive Operations}
          \begin{itemize}
            \item \textbf{Create Custom ZIP Archive} \label{sec:custom_archive}
                  
                  Creates a ZIP archive from arbitrary files with key-based manifest for
                  non-standard data. Call the \texttt{create\_zip\_archive} subroutine with
                  the following arguments:
                  
                  \textbf{Input arguments:}
                  \begin{itemize}
                    \item \texttt{zip\_filename character(len=*)}: Output ZIP file name
                          
                    \item \texttt{keys character(len=*)}: Array of keys for manifest entries
                          
                    \item \texttt{filenames character(len=*)}: Array of file paths to include
                          
                    \item \texttt{ierr integer(int32)}: Error code
                  \end{itemize}
                  \begin{lstlisting}[language=Fortran, caption={Create Custom Archive}]
character(len=32) :: keys(3) = ['test_data_1', 'test_data_2', &
                                'test_data_3']
character(len=256) :: files(3) = ['test_data_1.bin', 'test_2.bin', &
                                    'test_3.bin']
integer(int32) :: ierr

call create_zip_archive("custom_data.zip", keys, files, ierr)
                  \end{lstlisting}
                  
            \item \textbf{Extract ZIP Archive}
                  
                  Extracts all files from a ZIP archive to the disk in the current directory
                  and returns key-value pairs from the manifest. Call the \texttt{extract\_zip\_archive}
                  subroutine with the following arguments:
                  
                  \textbf{Input arguments:}
                  \begin{itemize}
                    \item \texttt{zip\_filename character(len=*)}: Input ZIP file name
                          
                    \item \texttt{ierr integer(int32)}: Error code
                  \end{itemize}
                  \textbf{Output arguments:}
                  \begin{itemize}
                    \item \texttt{keys character(len=:), allocatable}: Array of keys from
                          manifest
                          
                    \item \texttt{filenames character(len=:), allocatable}: Array of extracted
                          filenames
                  \end{itemize}
                  \begin{lstlisting}[language=Fortran, caption={Extract Custom Archive}]
character(len=:), allocatable :: archive_keys(:), archive_files(:)
integer(int32) :: ierr

call extract_zip_archive("custom_data.zip", archive_keys, &
                            archive_files, ierr)
                  \end{lstlisting}
                  
                  \textbf{Complete Archive Workflow:} 
                  \begin{lstlisting}[language=Fortran, caption={Complete Archive Management Workflow}]
integer(int32) :: ierr

! 1. Save processed data
call save_tox_data("analysis_results.zip", ierr, &
gene_ids=gene_ids, gene_ids_file="genes.bin", &
expression=expr_data, expression_file="expression.bin", &
family_centroids=centroids, family_centroids_file="centroids.bin")

! 2. Later, read the data back
character(len=:), allocatable :: loaded_genes(:)
real(real64), allocatable :: loaded_expr(:,:), loaded_centroids(:,:)
call read_tox_data("analysis_results.zip", ierr, &
                    gene_ids=loaded_genes, &
                    expression=loaded_expr, &
                    family_centroids=loaded_centroids)

! 3. Create custom archive with additional results
character(len=32) :: custom_keys(2) = ['data_1', 'data_2']
character(len=256) :: custom_files(2) = ['data_1_v2.bin', &
                                        'data_2_v1.bin']
call create_zip_archive("full_analysis.zip", custom_keys, &
                        custom_files, ierr)

! 4. Extract and verify custom archive
character(len=:), allocatable :: extracted_keys(:), extracted_files(:)
call extract_zip_archive("full_analysis.zip", extracted_keys, &
                        extracted_files, ierr)
            \end{lstlisting}
            
            \textbf{Key Features of Archive Operations:}
            \begin{itemize}
                \item \textbf{Automatic Memory Management}: Arrays automatically allocated
                      to correct sizes
                      
                \item \textbf{Flexible Data Selection}: Save/read any combination of
                      data arrays
                      
                \item \textbf{Cross-Platform Compatibility}: Archives work across R,
                      Python, and Fortran implementations
                      
                \item \textbf{Manifest-Based}: Automatic tracking of file contents and
                      relationships
                      
                \item \textbf{Temporary File Handling}: Automatic cleanup of intermediate
                      files
                      
                \item \textbf{Error Recovery}: Comprehensive error checking and reporting
            \end{itemize}
            
            \textbf{Archive Structure:} 
            \begin{verbatim}
archive.zip
|- manifest.txt          (key=filename mappings)
|- genes.bin            (gene IDs, if provided)
|- expression.bin       (expression data, if provided)
|- mapping.bin          (gene to family, if provided)
|- families.bin         (family IDs, if provided)
|- centroids.bin        (family centroids, if provided)
|- shifts.bin           (shift vectors, if provided)
            \end{verbatim}
          \end{itemize}
          
    \item \textbf{Individual Array Save/Load Operations}
          \begin{itemize}
            \item \textbf{Save Gene IDs}
                  
                  Saves a gene IDs array to a binary file with automatic character length
                  preservation. Call the \texttt{save\_gene\_ids} subroutine with the following
                  arguments:
                  
                  \textbf{Input arguments:}
                  \begin{itemize}
                    \item \texttt{gene\_ids character(len=*)}: Array of gene identifiers
                          
                    \item \texttt{filename character(len=*)}: Output binary file name
                          
                    \item \texttt{ierr integer(int32)}: Error code
                  \end{itemize}
                  \begin{lstlisting}[language=Fortran, caption={Save Gene IDs Array}]
character(len=256) :: gene_ids(10000)
integer(int32) :: ierr

call save_gene_ids(gene_ids, 'gene_ids.bin', ierr)
                  \end{lstlisting}
                  
            \item \textbf{Load Gene IDs}
                  
                  Loads a gene IDs array from a binary file with automatic allocation and
                  character length detection. Call the \texttt{load\_gene\_ids} subroutine
                  with the following arguments:
                  
                  \textbf{Input arguments:}
                  \begin{itemize}
                    \item \texttt{filename character(len=*)}: Input binary file name
                          
                    \item \texttt{ierr integer(int32)}: Error code
                  \end{itemize}
                  \textbf{Output arguments:}
                  \begin{itemize}
                    \item \texttt{gene\_ids character(len=:), allocatable}: Loaded gene IDs
                          array
                  \end{itemize}
                  \begin{lstlisting}[language=Fortran, caption={Load Gene IDs Array}]
character(len=:), allocatable :: loaded_gene_ids(:)
integer(int32) :: ierr

call load_gene_ids(loaded_gene_ids, 'gene_ids.bin', ierr)
! loaded_gene_ids automatically allocated with correct size 
! and character length
                  \end{lstlisting}
                  
            \item \textbf{Save Expression Vectors}
                  
                  Saves an expression data matrix to a binary file. Call the \texttt{save\_expression\_vectors}
                  subroutine with the following arguments:
                  
                  \textbf{Input arguments:}
                  \begin{itemize}
                    \item \texttt{expression real(real64)}: 2D expression data matrix
                          
                    \item \texttt{filename character(len=*)}: Output binary file name
                          
                    \item \texttt{ierr integer(int32)}: Error code
                  \end{itemize}
                  \begin{lstlisting}[language=Fortran, caption={Save Expression Data}]
real(real64) :: expr_data(n_samples, n_genes)
integer(int32) :: ierr

call save_expression_vectors(expr_data, 'expression.bin', ierr)
                  \end{lstlisting}
                  
            \item \textbf{Load Expression Vectors}
                  
                  Loads an expression data matrix from a binary file with automatic allocation.
                  Call the \texttt{load\_expression\_vectors} subroutine with the following
                  arguments:
                  
                  \textbf{Input arguments:}
                  \begin{itemize}
                    \item \texttt{filename character(len=*)}: Input binary file name
                          
                    \item \texttt{ierr integer(int32)}: Error code
                  \end{itemize}
                  \textbf{Output arguments:}
                  \begin{itemize}
                    \item \texttt{expression real(real64), allocatable}: Loaded expression
                          data matrix
                  \end{itemize}
                  \begin{lstlisting}[language=Fortran, caption={Load Expression Data}]
real(real64), allocatable :: loaded_expr(:,:)
integer(int32) :: ierr

call load_expression_vectors(loaded_expr, 'expression.bin', ierr)
! loaded_expr automatically allocated as (samples x genes)
                  \end{lstlisting}
                  
            \item \textbf{Save Gene to Family Mapping}
                  
                  Saves a gene to family mapping array to a binary file. Call the
                  \texttt{save\_gene\_to\_family} subroutine with the following arguments:
                  
                  \textbf{Input arguments:}
                  \begin{itemize}
                    \item \texttt{gene\_to\_family integer(int32)}: Gene to family mapping
                          array
                          
                    \item \texttt{filename character(len=*)}: Output binary file name
                          
                    \item \texttt{ierr integer(int32)}: Error code
                  \end{itemize}
                  \begin{lstlisting}[language=Fortran, caption={Save Gene-Family Mapping}]
integer(int32) :: gene_to_fam(n_genes)
integer(int32) :: ierr

call save_gene_to_family(gene_to_fam, 'gene_to_fam.bin', ierr)
                  \end{lstlisting}
                  
            \item \textbf{Load Gene to Family Mapping}
                  
                  Loads a gene to family mapping array from a binary file with automatic
                  allocation. Call the \texttt{load\_gene\_to\_family} subroutine with the
                  following arguments:
                  
                  \textbf{Input arguments:}
                  \begin{itemize}
                    \item \texttt{filename character(len=*)}: Input binary file name
                          
                    \item \texttt{ierr integer(int32)}: Error code
                  \end{itemize}
                  \textbf{Output arguments:}
                  \begin{itemize}
                    \item \texttt{gene\_to\_family integer(int32), allocatable}: Loaded gene
                          to family mapping
                  \end{itemize}
                  \begin{lstlisting}[language=Fortran, caption={Load Gene-Family Mapping}]
integer(int32), allocatable :: loaded_gene_to_fam(:)
integer(int32) :: ierr

call load_gene_to_family(loaded_gene_to_fam, 'gene_to_fam.bin', ierr)
                  \end{lstlisting}
                  
            \item \textbf{Save Family IDs}
                  
                  Saves a family IDs array to a binary file with automatic character length
                  preservation. Call the \texttt{save\_family\_ids} subroutine with the following
                  arguments:
                  
                  \textbf{Input arguments:}
                  \begin{itemize}
                    \item \texttt{family\_ids character(len=*)}: Array of family identifiers
                          
                    \item \texttt{filename character(len=*)}: Output binary file name
                          
                    \item \texttt{ierr integer(int32)}: Error code
                  \end{itemize}
                  \begin{lstlisting}[language=Fortran, caption={Save Family IDs}]
character(len=256) :: family_ids(n_families)
integer(int32) :: ierr

call save_family_ids(family_ids, 'family_ids.bin', ierr)
                  \end{lstlisting}
                  
            \item \textbf{Load Family IDs}
                  
                  Loads a family IDs array from a binary file with automatic allocation and
                  character length detection. Call the \texttt{load\_family\_ids} subroutine
                  with the following arguments:
                  
                  \textbf{Input arguments:}
                  \begin{itemize}
                    \item \texttt{filename character(len=*)}: Input binary file name
                          
                    \item \texttt{ierr integer(int32)}: Error code
                  \end{itemize}
                  \textbf{Output arguments:}
                  \begin{itemize}
                    \item \texttt{family\_ids character(len=:), allocatable}: Loaded family
                          IDs array
                  \end{itemize}
                  \begin{lstlisting}[language=Fortran, caption={Load Family IDs}]
character(len=:), allocatable :: loaded_family_ids(:)
integer(int32) :: ierr

call load_family_ids(loaded_family_ids, 'family_ids.bin', ierr)
                  \end{lstlisting}
                  
            \item \textbf{Save Family Centroids}
                  
                  Saves a family centroids matrix to a binary file. Call the \texttt{save\_family\_centroids}
                  subroutine with the following arguments:
                  
                  \textbf{Input arguments:}
                  \begin{itemize}
                    \item \texttt{family\_centroids real(real64)}: 2D family centroids matrix
                          
                    \item \texttt{filename character(len=*)}: Output binary file name
                          
                    \item \texttt{ierr integer(int32)}: Error code
                  \end{itemize}
                  \begin{lstlisting}[language=Fortran, caption={Save Family Centroids}]
real(real64) :: centroids(67, 15512)
integer(int32) :: ierr

call save_family_centroids(centroids, 'centroids.bin', ierr)
                  \end{lstlisting}
                  
            \item \textbf{Load Family Centroids}
                  
                  Loads a family centroids matrix from a binary file with automatic allocation.
                  Call the \texttt{load\_family\_centroids} subroutine with the following
                  arguments:
                  
                  \textbf{Input arguments:}
                  \begin{itemize}
                    \item \texttt{filename character(len=*)}: Input binary file name
                          
                    \item \texttt{ierr integer(int32)}: Error code
                  \end{itemize}
                  \textbf{Output arguments:}
                  \begin{itemize}
                    \item \texttt{family\_centroids real(real64), allocatable}: Loaded family
                          centroids matrix
                  \end{itemize}
                  \begin{lstlisting}[language=Fortran, caption={Load Family Centroids}]
real(real64), allocatable :: loaded_centroids(:,:)
integer(int32) :: ierr

call load_family_centroids(loaded_centroids, 'centroids.bin', ierr)
                  \end{lstlisting}
                  
            \item \textbf{Save Shift Vectors}
                  
                  Saves a shift vectors matrix to a binary file. Call the \texttt{save\_shift\_vectors}
                  subroutine with the following arguments:
                  
                  \textbf{Input arguments:}
                  \begin{itemize}
                    \item \texttt{shift\_vectors real(real64)}: 2D shift vectors matrix
                          
                    \item \texttt{filename character(len=*)}: Output binary file name
                          
                    \item \texttt{ierr integer(int32)}: Error code
                  \end{itemize}
                  \begin{lstlisting}[language=Fortran, caption={Save Shift Vectors}]
real(real64) :: shift_vecs(134, 88327)  ! 2xsamples x genes
integer(int32) :: ierr

call save_shift_vectors(shift_vecs, 'shift_vectors.bin', ierr)
                  \end{lstlisting}
                  
            \item \textbf{Load Shift Vectors}
                  
                  Loads a shift vectors matrix from a binary file with automatic allocation.
                  Call the \texttt{load\_shift\_vectors} subroutine with the following arguments:
                  
                  \textbf{Input arguments:}
                  \begin{itemize}
                    \item \texttt{filename character(len=*)}: Input binary file name
                          
                    \item \texttt{ierr integer(int32)}: Error code
                  \end{itemize}
                  \textbf{Output arguments:}
                  \begin{itemize}
                    \item \texttt{shift\_vectors real(real64), allocatable}: Loaded shift
                          vectors matrix
                  \end{itemize}
                  \begin{lstlisting}[language=Fortran, caption={Load Shift Vectors}]
real(real64), allocatable :: loaded_shifts(:,:)
integer(int32) :: ierr

call load_shift_vectors(loaded_shifts, 'shift_vectors.bin', ierr)
                  \end{lstlisting}
                  
            \item \textbf{Get Array Metadata}
                  
                  Retrieves metadata (dimensions, character length) from a binary file
                  without loading the entire array. Call the \texttt{get\_array\_metadata}
                  subroutine with the following arguments:
                  
                  \textbf{Input arguments:}
                  \begin{itemize}
                    \item \texttt{filename character(len=*)}: Binary file name to inspect
                          
                    \item \texttt{max\_dims integer(int32)}: Maximum number of dimensions
                          to return
                          
                    \item \texttt{ierr integer(int32)}: Error code
                  \end{itemize}
                  \textbf{Output arguments:}
                  \begin{itemize}
                    \item \texttt{dims integer(int32)}: Array dimensions
                          
                    \item \texttt{ndims integer(int32)}: Number of dimensions
                          
                    \item \texttt{char\_len integer(int32), optional}: Character length for
                          string arrays
                  \end{itemize}
                  \begin{lstlisting}[language=Fortran, caption={Inspect Array File Metadata}]
integer(int32) :: dims(5), ndims, char_len, ierr

! Get metadata for gene IDs file
call get_array_metadata('gene_ids.bin', dims, 5, ndims, &
                        ierr, char_len)

if (is_ok(ierr)) then
    print *, "Dimensions:", dims(1:ndims)
    print *, "Character length:", char_len
end if

! Get metadata for expression data file
call get_array_metadata('expression.bin', dims, 5, ndims, ierr)
if (is_ok(ierr)) then
    print *, "Expression data shape:", dims(1), "x", dims(2)
end if
                  \end{lstlisting}
                  
                  \textbf{Complete Individual Array Workflow:}
                  \begin{lstlisting}[language=Fortran, caption={Individual Array Save/Load Workflow}]
integer(int32) :: ierr

! 1. Save individual arrays
call save_gene_ids(gene_ids, 'gene_ids.bin', ierr)
call save_expression_vectors(expr_data, 'expression.bin', ierr)
call save_gene_to_family(gene_to_fam, 'mapping.bin', ierr)
call save_family_ids(family_ids, 'families.bin', ierr)
call save_family_centroids(centroids, 'centroids.bin', ierr)
call save_shift_vectors(shift_vecs, 'shifts.bin', ierr)

! 2. Check file metadata before loading
integer(int32) :: dims(5), ndims, char_len
call get_array_metadata('expression.bin', dims, 5, ndims, ierr)

! 3. Load arrays (automatic allocation)
character(len=:), allocatable :: loaded_genes(:)
real(real64), allocatable :: loaded_expr(:,:)
integer(int32), allocatable :: loaded_mapping(:)
character(len=:), allocatable :: loaded_families(:)
real(real64), allocatable :: loaded_centroids(:,:)
real(real64), allocatable :: loaded_shifts(:,:)

call load_gene_ids(loaded_genes, 'gene_ids.bin', ierr)
call load_expression_vectors(loaded_expr, 'expression.bin', ierr)
call load_gene_to_family(loaded_mapping, 'mapping.bin', ierr)
call load_family_ids(loaded_families, 'families.bin', ierr)
call load_family_centroids(loaded_centroids, 'centroids.bin', ierr)
call load_shift_vectors(loaded_shifts, 'shifts.bin', ierr)
                  \end{lstlisting}
                  
                  \textbf{Manual Archive Creation with Individual Arrays:}
                  \begin{lstlisting}[language=Fortran, caption={Manual Archive Creation}]
integer(int32) :: ierr

! 1. Save arrays to individual files
call save_gene_ids(gene_ids, 'temp_genes.bin', ierr)
call save_expression_vectors(expr_data, 'temp_expr.bin', ierr)

! 2. Create custom archive with specific keys
character(len=32) :: keys(2) = ['gene_data', 'expression_data']
character(len=256) :: files(2) = ['temp_genes.bin', 'temp_expr.bin']
call create_zip_archive('manual_archive.zip', keys, files, ierr)

! 3. Clean up temporary files (optional)
call delete_file('temp_genes.bin', ierr)
call delete_file('temp_expr.bin', ierr)
                  \end{lstlisting}
                  
                  \textbf{Key Features of Individual Array Operations:}
                  \begin{itemize}
                    \item \textbf{Automatic Allocation}: Output arrays automatically allocated
                          to correct sizes
                          
                    \item \textbf{Character Length Preservation}: String arrays maintain
                          original character lengths
                          
                    \item \textbf{Metadata Inspection}: Check array properties without loading
                          data
                          
                    \item \textbf{Binary Format}: Efficient storage for large datasets
                          
                    \item \textbf{Error Handling}: Comprehensive error checking for file
                          operations
                          
                    \item \textbf{Memory Efficient}: Can process arrays larger than available
                          memory
                  \end{itemize}
                  
                  \textbf{Performance Considerations:}
                  \begin{itemize}
                    \item Binary format provides fast I/O for large arrays
                          
                    \item Automatic dimension detection eliminates manual size tracking
                          
                    \item Character length preservation avoids fixed-width string limitations
                          
                    \item Suitable for datasets with millions of genes and thousands of samples
                  \end{itemize}
          \end{itemize}
\end{enumerate}